\begin{lemma}
\label{lemma-weighting-quasi-finite-Noetherian}
Let $f : X \to Y$ be a morphism of schemes. Assume
\begin{enumerate}
\item $f$ is locally quasi-finite, and
\item $Y$ is geometrically unibranch and locally Noetherian.
\end{enumerate}
Then there is a weighting $w : X \to \mathbf{Z}_{\geq 0}$ given by
the rule that sends $x \in X$ lying over $y \in Y$ to the
``generic separable degree''
of $\mathcal{O}_{X, x}^{sh}$ over $\mathcal{O}_{Y, y}^{sh}$.
\end{lemma}

\begin{proof}
It follows from Algebra, Lemma
\ref{algebra-lemma-quasi-finite-strict-henselization}
that $\mathcal{O}_{Y, y}^{sh} \to \mathcal{O}_{X, x}^{sh}$
is finite. Since $Y$ is geometrically unibranch
there is a unique minimal prime $\mathfrak p$ in
$\mathcal{O}_{Y, y}^{sh}$, see
More on Algebra, Lemma \ref{more-algebra-lemma-geometrically-unibranch}.
Write
$$
(\kappa(\mathfrak p) \otimes_{\mathcal{O}_{Y, y}^{sh}}
\mathcal{O}_{X, x}^{sh})_{red} =
\prod K_i
$$
as a finite product of fields.
We set $w(x) = \sum [K_i : \kappa(\mathfrak p)]_s$.

\medskip\noindent
Since this definition is clearly insensitive to \'etale localization,
in order to show that $w$ is a weighting we reduce to showing that if
$f$ is a finite morphism, then $\int_f w$ is locally constant.
Observe that the value of $\int_f w$ in a generic point $\eta$
of $Y$ is just the number of points of the geometric fibre
$X_{\overline{\eta}}$ of $X \to Y$ over $\eta$. Moreover, since
$Y$ is unibranch a point $y$ of $Y$ is the specialization of a unique
generic point $\eta$. Hence it suffices to show that $(\int_f w)(y)$
is equal to the number of points of $X_{\overline{\eta}}$.
After passing to an affine neighbourhood of $y$ we may assume
$X \to Y$ is given by a finite ring map $A \to B$. Suppose
$\mathcal{O}_{Y, y}^{sh}$ is constructed using a map
$\kappa(y) \to k$ into an algebraically closed field $k$.
Then
$$
\mathcal{O}_{Y, y}^{sh} \otimes_A B =
\prod\nolimits_{f(x) = y}
\prod\nolimits_{\varphi \in \Mor_{\kappa(y)}(\kappa(x), k)}
\mathcal{O}_{X, x}^{sh}
$$
by Algebra, Lemma \ref{algebra-lemma-finite-over-henselian}
and the lemma used above.
Observe that the minimal prime $\mathfrak p$ of $\mathcal{O}_{Y, y}^{sh}$
maps to the prime of $A$ corresponding to $\eta$. Hence we see that
the desired equality holds because the number of points of a geometric
fibre is unchanged by a field extension.
\end{proof}